\chapter{Discussion}
\label{ch:discussion}
 
\section{From Findings to Framework}
\label{sec:disc_intro}

The empirical evidence assembled in Chapter~\ref{ch:findings} is inconsistent with several assumptions that have organized post-\textit{Dobbs} legislative commentary. The modal 2023 abortion bill was neither a wholesale copy of an out-of-state template nor an act of local invention; across 188 enacted bills in 46 states, the mean cross-state novelty score was 0.769 ($\text{SD} = 0.122$), with two-thirds of bills clustered between 0.70 and 0.90. This concentration is the textual signature of what \textcite{karchDemocraticLaboratories2007} terms \emph{customization}: the selective incorporation of external frameworks with state-specific modifications. The first full legislative year after the constitutional vacuum opened was, in statistical terms, a year of adaptation, not of ex nihilo drafting.

Yet within that moderate-borrowing baseline, three results reverse the predicted direction and together demand theoretical reconstruction. First, it was protective coalitions, not restrictive ones, that deployed standardized templates most aggressively in 2023: protective bills are significantly less novel than neutral bills ($\beta = -0.076$, $p = 0.022$), while restrictive bills are statistically indistinguishable from the neutral baseline ($\beta = -0.017$, $p = 0.482$). Second, legislative professionalism is associated with \emph{greater} borrowing rather than innovation ($\beta = -0.155$, $p = 0.069$, marginal under Bell--McCaffrey CR2 small-sample corrections), directly reversing the intuition that well-staffed legislatures should produce more original legislation. Third, and most consequentially for the mechanism story, high textual coordination across state lines coexists with a near-complete absence of direct financial ties between coordinated legislators: zero common PAC donors appear across any of the ten highest-similarity cross-state bill pairs, including the Nevada SB~370--Washington HB~1155 pair that registers 80.7 percent cosine similarity, the singular case of near-replication in the corpus. The dyadic regression confirms at the state-pair level what the bill-level analysis suggested: shared protective coalition membership predicts significantly higher textual similarity ($\beta = +0.037$, $p < 0.001$), shared organizational network presence measured through PAC donor Jaccard overlap is an independent and comparable predictor ($\beta = +0.039$, $p < 0.001$), and geographic contiguity is null ($\beta = +0.003$, $p = 0.734$).

These findings resist integration into any single received account of policy diffusion. The geographic learning model of \textcite{walkerDiffusionInnovationsAmerican1969}, which assumes that neighbors influence neighbors, cannot explain why Nevada and Washington produced the strongest textual match in the dataset, or why Vermont serves as a template hub for Colorado, California, and New Mexico simultaneously. \textcite{berryStateLotteryAdoptions1990}'s event-history framework, which pairs internal determinants with regional diffusion, captures the internal determinants well but mistakes the diffusion channel. \textcite{jansaCopyPasteLawmaking2019}'s analysis of model-bill adoption correctly identifies interest-group drafting as a diffusion pathway but predicts, on the basis of pre-\textit{Dobbs} patterns, that the restrictive templating should dominate. None of these frameworks anticipated that the most visible cross-state coordination in 2023 would occur on the protective side, emerge from horizontal affiliate networks rather than vertical template distribution, and operate almost entirely outside the campaign-finance channel that most diffusion research has used as its proxy for interest-group influence.

This chapter develops the theoretical apparatus that the findings require. Section~\ref{sec:cvd} articulates \emph{constitutional vacuum diffusion} as a distinct diffusion regime triggered when federal preemption collapses, and sub-national jurisdictions are suddenly compelled to legislate across a newly opened policy frontier. Section~\ref{sec:rtm} develops \emph{reactive template mobilization} as the diffusion mechanism that operates under such a regime, distinguishing it from both learning-based diffusion and conventional crisis adoption. Section~\ref{sec:asymmetry} interprets the coalition asymmetry between horizontal and vertical diffusion architectures. Section~\ref{sec:reinterpreting-professionalism} reinterprets the professionalism finding as evidence of capacity for adoption rather than capacity for innovation. Section~\ref{sec:pac_paradox-discussion} develops the PAC paradox as empirical support for an informational rather than financial theory of interest-group influence. Section~\ref{sec:exporter-importer-reinforcer} revises the exporter--importer--reinforcer typology to accommodate the coalition-specific findings. Section~\ref{sec:scope-conditions} articulates the scope conditions under which the framework should travel. Section~\ref{sec:bridge} closes by indicating the inferential limits the observational design imposes on these claims and bridging to Chapter~\ref{ch:limitations}.

Table~\ref{tab:findings_crosswalk} summarizes the four empirical findings this chapter carries forward, each set against its conventional expectation and linked to the theoretical section in which it is developed. The remainder of the chapter proceeds through these findings in the order the table presents them, beginning with the diffusion regime within which the patterns emerge.

\begin{table}[!htbp]
\centering
\caption{Empirical Findings and Their Theoretical Destinations}
\label{tab:findings_crosswalk}
\renewcommand{\arraystretch}{1.4}
\setlength{\tabcolsep}{8pt}
\small
\begin{tabular}{>{\raggedright\arraybackslash}p{4.5cm} >{\raggedright\arraybackslash}p{3.8cm} >{\raggedright\arraybackslash}p{4.5cm}}
\toprule
\textbf{Empirical Finding (Ch.~\ref{ch:findings})} & \textbf{Conventional Expectation} & \textbf{Theoretical Destination (Ch.~\ref{ch:discussion})} \\
\midrule
Protective bills less novel than restrictive ($\beta = -0.076$, $p = 0.022$ vs.\ $\beta = -0.017$, $p = 0.482$) & Restrictive coalitions deploy templates more aggressively \parencite{jansaCopyPasteLawmaking2019} & Reactive template mobilization (\S\ref{sec:rtm}); coalition architecture asymmetry (\S\ref{sec:asymmetry}) \\
\addlinespace[6pt]
Professionalism inversely related to novelty ($\beta = -0.155$, $p = 0.069$) & Capacity enables independent innovation \parencite{squireMeasuringStateLegislative2007} & Capacity-for-adoption reinterpretation (\S\ref{sec:reinterpreting-professionalism}) \\
\addlinespace[6pt]
Zero PAC donor overlap across high-similarity pairs; Jaccard organizational overlap robust ($\beta = +0.039$, $p < 0.001$) & Financial contributions mediate interest-group influence & Informational logic of advocacy (\S\ref{sec:pac_paradox}) \\
\addlinespace[6pt]
Geographic contiguity null ($\beta = +0.003$, $p = 0.734$); Both Protective robust ($\beta = +0.037$, $p < 0.001$) & Regional diffusion \parencite{walkerDiffusionInnovationsAmerican1969,berryStateLotteryAdoptions1990} & Constitutional vacuum diffusion (\S\ref{sec:cvd}); ideological-network channel (\S\ref{sec:asymmetry}) \\
\bottomrule
\end{tabular}
\end{table}

\newpage
\section{Constitutional Vacuum Diffusion}
\label{sec:cvd}
The theoretical literature on policy diffusion has, with few exceptions, treated federalism as a background condition rather than as a variable. Whether diffusion is modeled as learning, imitation, competition, or coercion \parencite{shipanMechanismsPolicyDiffusion2008a}, the baseline assumption is that states operate within a stable constitutional allocation of authority and innovate at the margins of that allocation. What happens when the allocation itself shifts abruptly---when a domain that had been federally preempted for half a century is, without legislative warning, returned in full to the states---is a question the diffusion literature has had few occasions to answer.

\textit{Dobbs v. Jackson Women's Health Organization} (2022) is such an occasion. The decision did not merely modify the constitutional treatment of abortion; it vacated a half-century of federal doctrinal infrastructure within a single opinion. States that had spent decades legislating within the constraints of \textit{Roe} and \textit{Casey} were suddenly operating in a domain from which federal constitutional law had been withdrawn. I term the diffusion regime that emerges under such conditions \emph{constitutional vacuum diffusion}: the patterned interstate movement of policy text that occurs when federal preemption collapses, and states are compelled, within compressed legislative timelines, to either codify previously assumed protections or construct newly permissible restrictions.

Three features distinguish constitutional vacuum diffusion from the policy windows described by \textcite{kingdonAgendasAlternativesPublic2011a} and from the punctuated equilibrium dynamics theorized by \textcite{baumgartner_jones_1993}. First, the window is not discursively negotiated but judicially imposed: it opens when a court vacates prior doctrine rather than when a focusing event or problem stream aligns with political will. This matters for timing. \textcite{kingdonAgendasAlternativesPublic2011a}'s policy windows depend on the coincidence of three streams over which no single actor has full control, and the window may close before convergence. A judicially imposed vacuum opens on the day the opinion is issued and remains open until legislative action forecloses it, which creates a different temporal logic for actors on both sides of the policy debate.

Second, the vacuum is bilateral. Unlike conventional policy windows, which typically favor one coalition over another by opening space for a specific kind of reform, the constitutional vacuum creates urgent legislative imperatives for both coalitions simultaneously. States previously constrained by \textit{Roe}'s floor were free to legislate restrictions they could not previously enact; states that had relied on that same floor for protection were now compelled to codify in statute what constitutional doctrine had previously secured. The 88 restrictive and 73 protective bills enacted in 2023---a near-parity that would have been unthinkable under \textit{Roe}---reflect this dual imperative. Each coalition was legislating in response to the same event, but toward opposite ends and, crucially, from asymmetric starting positions.

Third, constitutional vacuum diffusion operates under conditions of acute informational asymmetry between coalitions. Restrictive coalitions entered 2023 with decades of accumulated statutory craft: gestational limits, waiting periods, informed consent requirements, ultrasound mandates, provider regulations, and a national model-bill infrastructure built by Americans United for Life, the Alliance Defending Freedom, and allied organizations specifically to test the boundaries of \textit{Roe} and \textit{Casey}. That craft was calibrated to the constraints of constitutional doctrine that had just been vacated. The total bans, trigger law activations, and enforcement mechanisms now permissible required statutory construction for which no litigation-tested template existed. Restrictive legislators in 2023 were, paradoxically, drafting language that registered no closer to standardized models than the neutral baseline, because the expansion of their policy frontier had outrun the template infrastructure their coalition had built.

Protective coalitions faced the opposite asymmetry. Their core policy instruments---shield laws, provider protection statutes, interstate cooperation provisions, emergency access codifications---had existed only in latent or marginal form before \textit{Dobbs}, because \textit{Roe} had rendered most of them unnecessary. Yet precisely because these frameworks had not been developed through decades of incremental statutory experimentation, they had to be \emph{constructed} rapidly in response to the vacuum. The organizational infrastructure that produced these frameworks---the Center for Reproductive Rights, the State Innovation Exchange, Planned Parenthood's state affiliate network, and NARAL---treated 2023 not as an occasion for novel drafting by individual state allies but as an occasion for the coordinated deployment of pre-positioned templates across every jurisdiction where the political arithmetic permitted enactment.

The constitutional vacuum diffusion framework therefore predicts a specific empirical signature: in the immediate aftermath of a preemption collapse, the coalition whose statutory toolkit was \emph{least} developed under the prior regime will exhibit the \emph{greatest} textual convergence, while the coalition whose toolkit was most developed will exhibit the greatest textual divergence. My findings match this prediction precisely. Protective coalitions inherited no refined template infrastructure and therefore deployed hastily coordinated standardized frameworks; restrictive coalitions inherited extensive template infrastructure calibrated to the prior constitutional regime and therefore found themselves drafting across a novel legal frontier for which their inherited templates were inadequate. The coalition asymmetry identified by H3 at the descriptive level is, under the constitutional vacuum diffusion framework, the predictable consequence of differential toolkit maturity at the moment the vacuum opened.

This framework also specifies the duration of the effect. Constitutional vacuum diffusion is, by construction, a transitional regime. As protective coalitions move beyond initial codification and as restrictive coalitions develop new litigation-tested templates for the expanded policy frontier, the asymmetry should attenuate. The 2023 textual signature captures a moment-specific equilibrium rather than a permanent structural feature, and the empirical patterns this thesis documents should not be expected to persist indefinitely. The framework's testable implication is that novelty scores for protective bills should rise over successive legislative cycles as the initial wave of shield law templating gives way to more idiosyncratic implementation questions, while novelty scores for restrictive bills should fall as new template infrastructure matures. Section~\ref{sec:scope-conditions} returns to this point as a scope condition. The 2023 textual signature is interpreted here against theoretical baselines drawn from the established diffusion literature; an explicit pre-\textit{Dobbs} corpus baseline, which would supply a sharper test of the framework's "departure" claim, is identified in Section~\ref{subsec:single-year} as a priority for longitudinal extension.

\section{Reactive Template Mobilization}
\label{sec:rtm}

Within the constitutional vacuum diffusion regime, the specific mechanism that produces the observed protective-coalition convergence is not adequately captured by any of the diffusion mechanisms previously theorized in the literature. This section develops \emph{reactive template mobilization} as a distinct mechanism, locates it within the existing taxonomy, and distinguishes it from adjacent concepts with which it should not be confused. 

Reactive template mobilization describes the rapid, coalition-coordinated deployment of pre-positioned legislative templates in direct response to the sudden expansion of the opposing coalition's policy frontier. Three features differentiate it from established diffusion mechanisms. First, it is \emph{reactive} rather than proactive: the mobilization is triggered not by endogenous coalition strategy but by an exogenous shock to the coalition's operating environment. Second, it deploys \emph{pre-positioned} templates---frameworks drafted, refined, and rehearsed within organizational networks before the shock but not yet activated for widespread legislative deployment. Third, it operates through \emph{simultaneous} adoption across multiple jurisdictions rather than sequential learning from early adopters.

The first feature distinguishes reactive template mobilization from the learning-based diffusion theorized by \textcite{shipanMechanismsPolicyDiffusion2008a}. Learning-based diffusion presumes that states observe the policy outcomes of early adopters, evaluate their effectiveness, and adapt the most successful templates for local conditions. This logic governs a substantial portion of pre-\textit{Dobbs} state policymaking on education, welfare, and regulatory matters. It does not fit the 2023 abortion record. Protective states did not wait to observe whether California's shield law produced the intended protective effects before replicating its language; the Nevada--Washington pair at 80.7 percent similarity was enacted within months of analogous California and Illinois statutes, before any meaningful evaluation of downstream effectiveness could have occurred. The temporal structure of the diffusion precludes the learning mechanism.

The second feature distinguishes reactive template mobilization from classic crisis adoption, in which states adopt emergency frameworks under time pressure precisely because they lack pre-developed policy tools. Crisis adoption in the conventional sense---the diffusion of emergency management statutes after natural disasters, or the spread of pandemic-response legislation in early 2020---involves actors constructing responses on short timelines with limited prior preparation. Reactive template mobilization is different: the templates exist before the shock. The organizational infrastructure had produced them, circulated them among affiliates, and refined them through internal review and model-bill drafting processes. What the shock triggers is not the creation of new policy content but the activation of content already in organizational reserve. The distinction matters analytically as it locates the causal work not in the urgency of the moment but in the anticipatory capacity of the organizational network.

The third feature distinguishes reactive mobilization from the classic network-based diffusion theorized by \textcite{ballaInterstateProfessionalAssociations2001} and elaborated in the advocacy coalition framework of \textcite{sabatier_weible_2007}. Network diffusion anticipates sequential adoption along established coalitional ties, typically moving from innovation-leader states through intermediate adopters to laggards, with each adoption informed by the previous ones. Reactive template mobilization produces a flatter temporal structure: states enact substantively similar legislation on compressed and overlapping timelines, not because they are learning from one another, but because they are activating the same pre-positioned template in response to the same exogenous shock. The Vermont hub pattern documented in \S\ref{sec:coordination}---in which Colorado SB 188, California SB 345, and New Mexico SB 13 all exhibit their closest cross-state matches to Vermont SB 37---is suggestive of this flatter structure, with a common template source rather than a sequential adoption chain.

Reactive template mobilization, therefore, occupies a distinct position in the diffusion taxonomy. It shares with network-based diffusion the reliance on organized cross-state advocacy infrastructure; it shares with crisis adoption the compressed temporal structure; it shares with learning-based diffusion the interstate movement of concrete statutory language. But it is reducible to none of them. The mechanism's distinctive contribution is its specification of the \emph{organizational anticipation} that enables rapid convergent adoption: it is the existence of templates in organizational reserve prior to the shock, not the learning processes or emergency responses that might follow it, that explains the textual signature observed in the 2023 corpus.

The practical implication for diffusion scholarship is that model-bill infrastructure matters most when the policy frontier shifts abruptly. In stable periods, organizational template reserves operate in the background, supplying drafting language to allied legislators without producing the statewide convergent signatures that text-as-data methods can detect. In periods of preemption collapse, the same reserves produce detectable signatures because the organizational anticipation is concentrated within a compressed adoption window. The 2023 abortion record is, from this vantage, not an exceptional case but an unusually legible one: the convergence is detectable because the shock forced simultaneous deployment of templates that would otherwise have been activated sequentially and invisibly over years.

Three rival mechanisms---parallel independent drafting from common doctrinal sources, informal staff copying from publicly available analysis, and post-enactment convergence produced by drafting conventions---remain consistent with the observed signature and are not eliminated by the present design. Each is enumerated explicitly in Section~\ref{subsec:text-necessary} and bracketed for the qualitative extension specified in Section~\ref{subsec:future-qualitative}.

The descriptive four-mechanism classification developed in Section~\ref{sec:four_mechanisms} (coordinated, collaborative, internal-learning, and competitive diffusion) describes patterns observable in the textual and documentary record. Reactive template mobilization, as articulated here, is the theoretical mechanism that produces the coordinated and collaborative subtypes under constitutional vacuum conditions; the descriptive scheme catalogs what is observed, while the theoretical scheme specifies what is responsible.

Table~\ref{tab:rtm-comparison} locates reactive template mobilization against the three established diffusion mechanisms with which it is most likely to be confused. The comparison proceeds along five analytic dimensions: the nature of the trigger, whether templates exist prior to the trigger, the temporal structure of adoption, the causal locus of the mechanism, and the scholarship most closely associated with each approach. The distinguishing feature of reactive template mobilization is not any single dimension but the specific combination: an exogenous trigger (shared with crisis adoption), pre-positioned templates (shared with network diffusion), and a compressed simultaneous temporal structure (shared with neither). This combination is what gives the mechanism its analytical distinctiveness.

Three rival mechanisms---parallel independent drafting from common doctrinal sources, information staff copying from publicly available analysis, and drafting-conventions convergence---are not eliminated by the present design and are enumerated explicitly in Section~\ref{subsec:text-necessary}.

\begin{landscape}
\begin{table}
\centering
\caption{Reactive Template Mobilization Compared to Established Diffusion Mechanisms}
\label{tab:rtm-comparison}
\renewcommand{\arraystretch}{1.3}
\setlength{\tabcolsep}{8pt}
\small
\begin{tabular}{@{}>{\raggedright\arraybackslash}p{3.2cm} >{\raggedright\arraybackslash}p{4.2cm} >{\raggedright\arraybackslash}p{4.2cm} >{\raggedright\arraybackslash}p{4.2cm} >{\raggedright\arraybackslash}p{4.6cm}@{}}
\toprule
\textbf{Feature} & \textbf{Learning} & \textbf{Network} & \textbf{Crisis} & \textbf{Reactive Template Mobilization} \\
\midrule
Trigger & Endogenous evaluation of prior adoptions & Endogenous coalition strategy & Exogenous emergency & Exogenous preemption collapse \\
\addlinespace
Templates exist before trigger? & Yes; adapted from observed cases & Yes; developed within coalition & No; constructed under pressure & Yes; prepositioned in organizational reserve \\
\addlinespace
Temporal structure & Sequential & Sequential along ties & Compressed, improvised & Compressed, simultaneous \\
\addlinespace
Causal locus & Outcome observation & Coalition tie strength & Time pressure & Organizational anticipation \\
\addlinespace
Comparative advantage under shock conditions & Slow; requires observable outcomes from prior adopters & Moderate; bounded by tie density & Rapid but improvisational; output quality variable & Rapid and coordinated; output quality preserved through prior refinement \\
\addlinespace
Representative scholarship & \textcite{shipanMechanismsPolicyDiffusion2008a} & \textcite{ballaInterstateProfessionalAssociations2001}; \textcite{sabatier_weible_2007} & \textcite{bousheyPolicyDiffusionDynamics2010a} & This thesis \\
\bottomrule
\end{tabular}
\end{table}
\end{landscape}

\section{Coalition Asymmetry and the Architecture of Advocacy Networks}
\label{sec:asymmetry}

The empirical coalition asymmetry---protective convergence in the bill-level analysis ($\beta = -0.076$, $p = 0.022$), restrictive non-significance ($\beta = -0.017$, $p = 0.482$), and the parallel structure at the dyadic level where Both Protective is robust ($\beta = +0.037$, $p < 0.001$) while Both Restrictive is null in the preferred specification ($\beta = +0.012$, $p = 0.158$)---admits a specific organizational interpretation. The difference between protective and restrictive coalitions in 2023 was not primarily a difference in template sophistication, coordination intensity, or financial resources. It was a difference in the \emph{architecture} of the networks through which templates moved.

Restrictive coalitions operate, and have operated since the 1970s, through vertical template distribution. National organizations---Americans United for Life, the Alliance Defending Freedom, Susan B. Anthony Pro-Life America, and allied bodies---draft model legislation, maintain it in publicly or semi-publicly accessible libraries, and distribute it to state legislative allies through a combination of direct lobbying, legislative staff liaisons, and model-bill clearinghouse functions. The structure is hierarchical: a small number of national drafting nodes produce templates that diffuse downward and outward to a large number of state-level adopting nodes. Within this architecture, individual state bills tend to be either near-verbatim adoptions of national models or substantial state-specific departures from them. There is relatively little horizontal adaptation among adopters.

The attenuation of the Both Restrictive coefficient between dyadic Models~D2 and~D3 is diagnostic of this vertical architecture. In Model~D2, shared restrictive coalition membership registers at the margin of significance ($\beta = +0.015$, $p < 0.10$); once the Jaccard index of shared PAC donor organizations enters Model~D3, the coefficient attenuates to null ($\beta = +0.012$, $p = 0.158$). The interpretation is straightforward: the marginal signal that raw restrictive coalition co-membership carried in D2 was absorbed by the organizational overlap variable in D3. That is, the mechanism by which two restrictive states come to enact textually similar legislation is substantially captured by the fact that the same national organizations maintain an advocacy footprint in both states. Coalition membership without organizational infrastructure does not produce convergence; organizational infrastructure produces it, and coalition co-membership is simply a proxy for that infrastructure.

Protective coalitions operate through a different architecture: horizontal coordination among state affiliates. The Center for Reproductive Rights, the State Innovation Exchange, Planned Parenthood's affiliate network, and NARAL do maintain national drafting capacity, but their state affiliates do not function primarily as passive recipients of national models. They function as co-drafters within distributed working groups that share frameworks laterally, adapt each other's statutory language, and refine common templates through iterative revision across jurisdictions. Within this architecture, the California–New York–Washington–Illinois cluster operates not as a center with peripheries, but as a set of mutually influential nodes, each of which originates some content and incorporates some content from the others.

The invariance of the Both Protective coefficient across dyadic Models~D2 and~D3---stable at $\beta = +0.037$ regardless of whether PAC organizational overlap is in the model---is diagnostic of this horizontal architecture. Unlike the restrictive case, the protective coalition effect is not substantially absorbed by the Jaccard variable. The interpretation is that protective convergence operates through channels only partially captured by shared PAC donor presence: affiliate working groups, cross-state legislative staff networks, coordinated testimony circulation, and shared legal-academic drafting capacity are all mechanisms of horizontal coordination that the Jaccard index does not measure. The measure captures organizational presence but not the distributed collaborative processes through which protective templates are actually constructed.

This architectural distinction has implications that extend beyond the 2023 abortion record. Vertical diffusion architectures are, in the language of network analysis, star-shaped: they produce high convergence when active, but they are vulnerable to the removal of central nodes, and they generate relatively homogeneous content because adaptation is penalized at the state level. Horizontal diffusion architectures are mesh-shaped: they exhibit moderate convergence during stable periods, yet they are resilient to node removal and capable of rapid scaling because any affiliate can initiate and any can adopt. Under the conditions of constitutional vacuum diffusion, the mesh architecture demonstrated comparative advantages the literature has not previously documented: it mobilized faster, achieved wider coverage, and sustained convergence across a broader set of dyads than the star architecture could match within the same legislative cycle. Whether this comparative advantage persists beyond the initial post-\textit{Dobbs} window or is specific to the preemption-collapse regime is an empirical question the current design cannot resolve, but Section~\ref{subsec:future-longitudinal} flags as a priority for extension.

\begin{table}[H]
\centering
\caption{Vertical and Horizontal Coalition Architectures in 2023 Abortion Diffusion}
\label{tab:architecture}
\renewcommand{\arraystretch}{1.4}
\setlength{\tabcolsep}{8pt}
\small
\begin{tabular}{>{\raggedright\arraybackslash}p{3.4cm} >{\raggedright\arraybackslash}p{4.6cm} >{\raggedright\arraybackslash}p{4.6cm}}
\toprule
\textbf{Dimension} & \textbf{Vertical (Restrictive)} & \textbf{Horizontal (Protective)} \\
\midrule
Network topology & Star-shaped: national drafting nodes to state adopters & Mesh-shaped: distributed affiliate co-drafting \\
\addlinespace[6pt]
Anchor organizations & Americans United for Life, Alliance Defending Freedom, Susan B.\ Anthony Pro-Life America & Center for Reproductive Rights, State Innovation Exchange, Planned Parenthood affiliates, NARAL \\
\addlinespace[6pt]
Drafting locus & National model-bill libraries & Cross-state working groups \\
\addlinespace[6pt]
Adaptation & Minimal at adopter level & Iterative across affiliates \\
\addlinespace[6pt]
Dyadic signature & Both Restrictive: $\beta = +0.012$, $p = 0.158$ (D3); attenuated by PAC Jaccard between D2 and D3 & Both Protective: $\beta = +0.037$, $p < 0.001$ (D3); invariant across D2 and D3 \\
\addlinespace[6pt]
Diagnostic interpretation & Convergence operates through shared organizational footprint; Jaccard absorbs the signal & Convergence operates through channels only partially captured by Jaccard (affiliate networks, staff coordination) \\
\addlinespace[6pt]
Comparative advantage under shock & Slower mobilization; constrained by central node bottleneck & Faster mobilization; broader coverage; resilient to node loss \\
\bottomrule
\end{tabular}
\end{table}

\section{Reinterpreting Legislative Professionalism}
\label{sec:reinterpreting-professionalism}

The directional inversion of H1---a negative rather than positive relationship between the Squire Index and cross-state novelty, with $\beta = -0.155$ ($p = 0.069$) in the full model---deserves separate theoretical treatment because it revises a long-standing intuition in the state politics literature. The received view, developed across \textcite{squireLegislativeProfessionalizationMembership1992a}, \textcite{squireMeasuringStateLegislative2007}, and the body of work applying these measures to state policy output, has been that professionalism enhances legislative capacity, that capacity enables independent policy development, and that independent policy development manifests as higher rates of original legislation. The 2023 findings suggest the causal chain requires a different specification.

The most parsimonious reinterpretation, consistent with \textcite{jansaCopyPasteLawmaking2019} on model-bill adoption and with \textcite{shipanMechanismsPolicyDiffusion2008a} on the conditional operation of diffusion mechanisms, is that legislative professionalism enables \emph{adoption} rather than \emph{innovation}. Professional legislatures possess the research staff, the committee capacity, and the institutional memory to identify external templates, evaluate their legal and political viability, and strategically adapt them to local political and institutional contexts. Less-professionalized legislatures, by contrast, lack the systematic infrastructure to identify, evaluate, and modify external templates. 

Under this interpretation, templates are produced primarily by national advocacy organizations and affiliated research arms whose drafting capacity operates outside the state legislative system. Americans United for Life, the Alliance Defending Freedom, the Center for Reproductive Rights, and the State Innovation Exchange maintain full-time legal and policy staff whose function is precisely to construct and refine model statutory language for legislative allies. Professional state legislatures have the staff capacity to locate, evaluate, and adapt this externally produced material; low-professionalism legislatures do not, and the textual originality they exhibit is therefore a signature of informational isolation from the template infrastructure rather than evidence of independent innovation. The causal work, in other words, is done not by the adopting legislatures but by the organizational drafting nodes that sit upstream of them---a distinction the bill-level analysis in Chapter~4 cannot directly observe, and which the current design establishes only by inference from the pattern of convergence across coalitionally aligned states.

This does not mean that low-professionalism legislatures produce truly original legislation in the substantive sense; it means they draft more independently because they lack access to the template infrastructure that professional legislatures can exploit.

The textual originality that low-professionalism states exhibit is, on this reading, a signature of informational isolation rather than creative capacity. Wyoming's high novelty (mean 0.82, Squire Index 0.071) and Maine's comparable score (mean 0.85, Squire Index 0.133) do not indicate that these legislatures are producing more innovative policy than Michigan (mean 0.86, Squire Index 0.657); they indicate that these legislatures are producing less templated policy because their access to template infrastructure is more limited.

This reinterpretation yields a testable implication: the negative professionalism--novelty relationship should be strongest in policy domains where organizational template infrastructure is dense, and should weaken or reverse in domains where such infrastructure is sparse. Abortion policy in 2023 is a high-infrastructure case: decades of model-bill production by both coalitions mean that any state with the staff capacity to access those templates can deploy them, while states without that capacity cannot. In a low-infrastructure domain—a newly emergent regulatory area with no established advocacy organizations or model-bill clearinghouses—professionalism should manifest as innovation rather than adoption, because there is no external template to adopt. The current design does not test this cross-domain prediction, but it supplies the mechanism that such a test would require.

The further implication for diffusion scholarship is that text-based measures of novelty should not be read as measures of policy innovation in the normative sense. A state that adopts an externally drafted bill well-calibrated to its local conditions may be engaging in more sophisticated policy development than one that drafts a bill from scratch. The novelty score captures the textual signature of that adoption, but the evaluative inference---whether the adoption reflects good or poor policymaking---requires criteria the text-based measure cannot supply. Chapter~\ref{ch:limitations} returns to this point as a methodological caveat.

Direct empirical purchase on this interpretation would require evidence that the current design cannot supply: observation of the drafting processes internal to legislative offices, including the channels through which staff in high- and low-professionalism states access or fail to access externally produced templates. Chapter~6 identifies this form of inquiry---qualitative investigation of legislative drafting processes---as the most promising extension of the research program initiated here, and as the most direct test of whether the capacity-for-adoption reinterpretation survives closer examination of the drafting environment itself.

\section{The PAC Paradox and the Informational Logic of Advocacy Networks}
\label{sec:pac_paradox-discussion}

The PAC paradox---high textual coordination coexisting with near-zero direct financial connections between coordinated legislators---constitutes perhaps the most consequential finding for the broader interest-group and policy-diffusion literature. Cross-referencing the ten highest-similarity cross-state pairs with Federal Election Commission data yielded zero common PAC donors. The Nevada--Washington pair, the strongest coordination case in the corpus, involved sponsors neither of whom received any abortion-related PAC contributions in the 2022 cycle. Systematic analysis of the DIME database's full 53 million contribution records confirmed near-zero direct contributions from abortion advocacy organizations to state legislative candidates. Yet the dyadic regression shows that the Jaccard index of shared PAC donor \emph{organizations} between two states is a robust predictor of textual similarity ($\beta = +0.039$, $p < 0.001$).

The reconciliation is straightforward but has substantial implications. The PAC measure that predicts convergence is not the contribution itself but the organizational presence it signals. Two states in which Planned Parenthood, Susan B. Anthony Pro-Life America, NARAL, and the National Right to Life Committee maintain concurrent advocacy footprints exhibit greater textual similarity in their abortion legislation, not because those organizations have purchased legislative behavior in either state, but because their concurrent presence indicates a shared informational and organizational ecosystem. The organizations supply legislative drafts, legal memoranda, model statutory language, expert testimony, coalition talking points, and staff liaison services; they do not meaningfully provide financial support to state legislative campaigns. The convergence the dyadic regression detects is the textual imprint of that informational coordination, not of any financial transaction.

The mechanical question of what activity places these organizations in the DIME database---and therefore into the Jaccard index---warrants direct clarification, since the index would be uninformative if the relevant organizations were absent from the data altogether. DIME records contributions from organized donors to any candidate, committee, or party apparatus tracked by the Federal Election Commission and its state-level analogues, which encompasses a substantially broader set of transactions than contributions to individual state legislative candidates. Planned Parenthood, Susan B. Anthony Pro-Life America, NARAL, and the National Right to Life Committee appear in DIME through contributions to gubernatorial campaigns, state party committees, ballot-initiative organizations, federal officeholders representing the state, and affiliated political action committees that themselves operate at the state or local level. The Jaccard index therefore captures the \textit{co-presence} of these organizations within a state's broader political economy, registering whether both states host the same advocacy infrastructure in any electorally meaningful form. It does not capture, and is not intended to capture, the magnitude of financial transfers to the specific legislators who sponsored the bills analyzed in this thesis. That latter quantity is precisely what the analysis in Section~4.8 shows to be negligible, and the decoupling between the two measures is exactly the empirical pattern the informational-subsidy interpretation predicts: the organizations whose presence predicts textual convergence are the same organizations whose direct contributions to bill sponsors approach zero.

This finding aligns with and extends the theoretical framework advanced in \textcite{hallLobbyingLegislativeSubsidy2006a}. Hall and Deardorff re-conceptualize interest-group influence as legislative subsidy: organizations affect policy not by purchasing favorable votes but by reducing the cost of favorable legislative action for allies who would already have supported the position on ideological grounds. The subsidy takes the form of research, drafting, constituent communications, and procedural assistance. \textcite{hertel-fernandezStateCapture2019} documents a parallel dynamic at the state level, showing that organizations like Americans United for Life and Susan B. Anthony Pro-Life America operate primarily through model-bill distribution, policy conferences, and staff liaison rather than through direct contribution channels. The 2023 abortion findings provide empirical confirmation of both arguments in a policy domain where the financial-influence intuition has been especially strong in public commentary.

The implication for diffusion research methodology is that campaign finance data, used alone, are an inadequate proxy for interest-group influence on state legislation. A study that measured advocacy effects solely through contribution flows would have concluded, from the 2023 abortion record, that abortion advocacy organizations exert minimal influence on state legislative text. The textual evidence demonstrates the opposite. The contribution flows are small because the mechanism does not run through them; the convergence is substantial because it runs through channels that contribution data do not capture. Future diffusion research that aims to measure organizational influence should combine financial measures with textual, network, or organizational-presence measures---and should be prepared to find the financial measures uninformative or misleading in domains where the informational channel dominates.
 
A secondary implication concerns how the literature has theorized the distinction between financial and informational mechanisms. In \textcite{hallLobbyingLegislativeSubsidy2006a}'s original formulation, the two mechanisms are presented as complementary: organizations subsidize the legislators they also contribute to. The 2023 abortion findings suggest the mechanisms can operate with substantial independence. Organizations can exert significant influence on legislators to whom they contribute nothing, provided the ideological alignment is strong and the template infrastructure is well developed. This decoupling of the financial and informational channels may be characteristic of polarized policy domains where ideological sorting is sufficiently complete that contributions add little to what ideological alignment already supplies. Testing this conjecture across other polarized domains---gun policy, environmental regulation, voting rights---is a natural extension of the research program this thesis initiates.

A methodological caveat is in order before these claims travel further. The interpretation developed in this section rests on indirect evidence. The current design establishes, with considerable precision, that textual convergence across coalitionally aligned states coexists with the presence of shared advocacy organizations in those states' political economies, and that this convergence coexists with the near-complete absence of direct financial ties between the coordinated legislators themselves. What the design does not directly observe is the drafting process through which model language actually moves from organizational nodes into enacted legislation. The informational-subsidy interpretation is the most parsimonious account of the observed pattern and is consistent with the broader literature on organizational influence at the state level \parencite{hertel-fernandezStateCapture2019,hallLobbyingLegislativeSubsidy2006a}, but it remains an inference from the structure of convergence rather than a direct observation of the mechanism. Chapter~6 returns to this limitation at length and identifies the forms of evidence---chiefly direct inquiry into the drafting processes internal to legislative offices and advocacy organizations---that would be required to strengthen the mechanism-level claim beyond what the current observational design can support.
 
\section{Revisiting the Exporter--Importer--Reinforcer Typology}
\label{sec:exporter-importer-reinforcer}
 
The exporter--importer--reinforcer typology developed in Chapter~\ref{ch:theory} anticipated that high-capacity states would produce original legislation subsequently adopted elsewhere (exporters), that lower-capacity states would adopt external models (importers), and that a third category would rework existing frameworks with minimal substantive deviation (reinforcers). The 2023 findings both support and complicate this typology, and the complications are theoretically productive rather than disconfirming.
 
The core insight holds: states occupy distinct functional positions in the diffusion network, which can be identified through textual evidence and documentary corroboration. What requires revision are the assumptions that exporter status tracks institutional capacity and that the same typology applies uniformly across coalitions.
 
The professionalism finding refutes the first assumption. If California, with its Squire Index of 0.626, functions in 2023 as a net importer rather than as the expected exporter---incorporating framework language from New York, Washington, and other protective-coalition partners---then the equation between legislative capacity and exporter status does not hold. The state-level evidence in §4.3 makes this concrete. California's mean novelty of 0.75 sits slightly below the corpus mean of 0.769, and its closest cross-state matches trace to Vermont, New York, and other protective-coalition nodes rather than to states importing California language. California functions in the 2023 record as a net adapter rather than the expected net exporter—a position that reverses the conventional mapping between high legislative capacity (Squire = 0.626 in the 2021 update; \parencite{squireSquireIndexUpdate2024}) and exporter status. Exporter status is therefore not a function of capacity in isolation; it is a function of capacity interacting with coalition architecture and with the temporal position of the state in the post-shock adoption sequence.
 
The coalition asymmetry refutes the second assumption. Within the vertical restrictive architecture, the typology operates cleanly: a small number of national organizational nodes (not individual states) generate templates that a larger number of state legislatures import with minimal modification. Texas (mean novelty 0.83) functions as the clearest state-level exporter in this architecture, drafting frameworks subsequently incorporated by Mississippi, Louisiana, and Arkansas, though its exporter status is substantially mediated by the national organizations that have historically used Texas statutes as testbeds for model legislation. Utah occupies a related position: its 2023 enactments are too few in number to qualify it as an exporter at the state level, but Utah-origin templates from preceding legislative cycles—particularly the trigger-law architecture and definitional provisions developed in 2018–2020—appear in the 2023 corpus through downstream adoption in Arkansas, Montana, and Wyoming. Arkansas, Montana, Wyoming, and similar states function as importers of this template flow. The reinforcer category fits Virginia and Mississippi, which reworked existing definitional frameworks with minimal substantive deviation across multiple 2023 enactments.
 
Within the horizontal protective architecture, the typology operates differently, and the categories themselves become more fluid. California functions simultaneously as an exporter (originating the framework for AB~100 and related funding bills), an adapter (incorporating shield law language from Vermont and Washington), and a reinforcer (reworking its own prior codifications without substantive expansion). Vermont functions as a template hub, with Vermont SB~37 serving as the closest cross-state match for bills in Colorado, California, and New Mexico. Vermont's middling legislative capacity (Squire Index 0.301) makes the hub function theoretically interesting: it cannot be attributed to the staff resources and research infrastructure that explain why high-capacity states sometimes originate widely adopted frameworks, nor to the informational isolation that explains why low-capacity states sometimes appear original by default. The hub function is instead a product of coordinated template development by protective-coalition affiliates, which culminated in Vermont legislation early in the 2023 cycle.

The expectation that higher-professionalism states should move first in any diffusion sequence---an expectation consistent with \textcite{squireMeasuringStateLegislative2007}, \textcite{shipanMechanismsPolicyDiffusion2008a}, and the broader state politics literature on legislative capacity---is not borne out in the Vermont case. California, New York, and Washington each possess substantially greater legislative capacity than Vermont, and each enacted protective legislation in 2023, yet the textual evidence positions Vermont SB 37 as the closest cross-state match for bills in Colorado, California, and New Mexico rather than the reverse. Two interpretations are available, and they are not mutually exclusive. The first is that template hub status within the protective coalition is determined not by the individual capacity of adopting legislatures but by the position of affiliated advocacy organizations within the horizontal drafting network, with Vermont's early 2023 enactment reflecting coordinated affiliate timing rather than independent institutional capacity. The second is that high-capacity states, precisely because their legislative processes are slower, more formally institutionalized, and more procedurally encumbered, may move later in the adoption sequence than lower-capacity states whose smaller legislatures can enact protective legislation on compressed timelines. The current design cannot adjudicate between these interpretations, and both would be consistent with the horizontal architecture developed in Section~5.4. What the Vermont case does establish, regardless of which interpretation prevails, is that the conventional mapping between institutional capacity and diffusion function---in which high-capacity states originate frameworks and low-capacity states adopt them---requires substantial revision to accommodate the empirical record this thesis documents. Future research capable of modeling enactment timing explicitly, identified in Chapter~6 as a priority for longitudinal extension, would supply the additional traction this question requires.
 
The revised typology, accommodating these findings, has three features the original lacked. First, it distinguishes exporter \emph{status} from exporter \emph{function}: a state can originate a framework that diffuses widely (function) without maintaining high levels of original drafting across its legislative output (status). In 2023, Vermont occupied this position on protective legislation. Second, it acknowledges that the importer--exporter axis operates within a coalition structure rather than across one: comparing a high-novelty restrictive state to a low-novelty protective state confuses two distinct diffusion architectures that should be analyzed within-coalition before being compared across-coalition. Third, it treats the reinforcer category as genuinely distinct from the importer category on grounds of organizational intent rather than textual signature alone: a state that reworks existing language to maintain policy continuity (reinforcer) is doing different organizational work than a state that adopts external language to import new policy (importer), even when the textual novelty scores are comparable.
 
These refinements do not abandon the original typology; they specify it more carefully. The exporter--importer--reinforcer categories remain analytically useful descriptors of state function in diffusion networks. The contribution of the 2023 analysis is to show that the categories operate differently across coalition architectures and that institutional capacity, while relevant, does not map cleanly onto functional position in the ways the pre-\textit{Dobbs} literature on state innovation suggested.
 
\section{Scope Conditions}
\label{sec:scope-conditions}
 
The theoretical apparatus developed across the preceding sections is not intended to travel uniformly across all domains of state policy diffusion. This section articulates the scope conditions under which constitutional vacuum diffusion, reactive template mobilization, and the coalition-architecture distinction should be expected to apply. It identifies the conditions under which the pre-existing diffusion literature should retain its explanatory priority.
 
\subsection{Preemption Collapse Event}
The first and most consequential scope condition is the presence of a \emph{preemption collapse event}. Constitutional vacuum diffusion is, by definition, a regime that emerges when federal constitutional or statutory authority that had previously constrained state action is abruptly withdrawn. Such events are rare. The \textit{Dobbs} decision is the clearest example within American federalism in the last half-century. Analogous events are few: the partial invalidation of the Voting Rights Act's preclearance regime in \textcite{ShelbyCountyAla2013}, the striking of the individual mandate's enforcement mechanism in the Affordable Care Act litigation cycle, and the post-Heller and post-Bruen recalibration of permissible firearms regulation following \textcite{ussc2008heller} and \textcite{nyrifle2022bruen} are candidate cases, though each differs from \textit{Dobbs} in the completeness and suddenness of the doctrinal withdrawal. Outside of these rare events, the diffusion dynamics in most policy domains are governed by the incremental processes theorized by the established literature. The framework developed here is a framework for anomalous periods, not a replacement for the diffusion baseline.
 
\subsection{Bilateral Pre-Existing Advocacy Infrastructure}
The second scope condition is the \emph{presence of dense pre-existing advocacy infrastructure on both sides of the policy debate}. Reactive template mobilization requires that at least one coalition have pre-positioned templates available for rapid deployment. In policy domains where no such templates exist---newly emergent regulatory questions without established organizational presence---the shock of a preemption collapse would not produce the textual convergence the 2023 abortion record exhibits, because there would be nothing to deploy. The asymmetry between protective and restrictive coalitions in 2023 depended on both coalitions having \emph{some} organizational infrastructure; had protective coalitions lacked the shield law frameworks that the Center for Reproductive Rights and allied organizations had been developing, the post-\textit{Dobbs} legislative record would likely have exhibited the scattered, independent-drafting signature that low-capacity innovation produces. The framework, therefore, applies most cleanly to policy domains with long histories of organized advocacy on both sides.
 
\subsection{Temporal Proximity to the Shock}
The third scope condition is \emph{temporal proximity to the shock}. The 2023 findings capture legislative behavior during the first full calendar year after \textit{Dobbs}. The theoretical argument in Section~\ref{sec:cvd} predicts that the coalition asymmetry should attenuate as the protective coalition's initial codification wave subsides and the restrictive coalition develops new litigation-tested templates for the expanded policy frontier. The framework's applicability to 2024, 2025, and subsequent legislative cycles is therefore an empirical question the current design cannot resolve. The framework provides a testable prediction about the direction of change: protective novelty should rise, and restrictive novelty should fall as the post-shock adjustment period matures. Chapter~\ref{ch:limitations} treats this prediction as the foremost candidate for future longitudinal extension.
 
\subsection{Decoupled Financial and Information Channels}
The fourth scope condition is the \emph{decoupling of financial and informational influence channels}. The PAC paradox documented in Section~\ref{sec:pac_paradox-discussion} may be characteristic of highly polarized, ideologically sorted policy domains in which legislators' positions are pre-determined by coalitional identity and organizational contributions add little to what ideology already supplies. In less polarized domains where legislators' votes are not pre-determined---infrastructure, taxation, many regulatory questions---the financial and informational channels may operate in closer coordination, and the PAC paradox would not be expected to appear. The framework's decoupling claim, therefore, applies most cleanly to polarized domains and should not be extended to the full range of state legislative activity without additional empirical testing.
 
These scope conditions are not limitations in the pejorative sense; they are the boundary specifications that a theoretical framework requires if it is to be useful for comparative work. Constitutional vacuum diffusion, reactive template mobilization, and the coalition-architecture distinction are not intended as general theories of state policymaking. They are specific mechanisms that operate under specific conditions, and their value for the broader diffusion literature depends on accurate specification of when they should and should not be expected to apply.
 
\section{Implications and Bridge to Limitations}
\label{sec:bridge}
 
The theoretical apparatus developed in this chapter makes several contributions that Chapter~7 will develop at greater length. It specifies a diffusion regime the pre-\textit{Dobbs} literature did not anticipate; it introduces a mechanism---reactive template mobilization---that occupies a distinct position in the existing taxonomy; it documents coalition architectures that produce differentiated convergence signatures under shock conditions; it reinterprets a long-standing assumption about legislative professionalism; and it provides evidence consistent with the informational-subsidy account of interest-group influence in a domain where the financial-influence intuition has dominated public commentary.
 
These contributions, however, rest on inferential foundations that the observational research design can support only partially. The quantitative evidence establishes textual convergence with considerable precision: the TF-IDF vectorization over 104,445 unique n-grams, the 188$\times$188 cosine similarity matrix, the Bell--McCaffrey CR2 cluster-robust standard errors computed via \texttt{estimatr::lm\_robust} following \textcite{pustejovskyCRVEAnywhere2018}, and the dyadic specification on 1,035 state pairs collectively provide a rigorous description of the textual patterns that characterized 2023 post-\textit{Dobbs} legislation. What the evidence cannot fully establish is the organizational mechanism responsible for producing those patterns. The documentary evidence examined in Section~\ref{sec:org_networks}---LegiScan legislative records, FEC contribution data, advocacy organization publications, and legislative history materials---corroborates the organizational interpretation, but documentary evidence speaks to organizational presence and public activity rather than to the internal drafting processes through which specific statutory language is developed and circulated.
 
The gap between the textual pattern and the organizational mechanism is the principal inferential limit imposed by the current design. The constitutional vacuum diffusion framework predicts that protective coalitions in 2023 deployed pre-positioned templates through horizontal affiliate coordination; the evidence is consistent with this prediction but does not fully establish it. Alternative mechanisms---parallel independent drafting by coalitionally aligned legislators working from common doctrinal sources, legislative staff copying from academic and advocacy publications, post-enactment textual convergence produced by drafting conventions rather than intentional coordination---are not eliminated by the current design. Distinguishing among these candidate mechanisms with confidence would require direct evidence about the drafting processes within legislative offices and advocacy organizations in 2023, evidence that the computational and documentary materials assembled for this thesis cannot provide.
 
Chapter~\ref{ch:limitations} addresses these inferential limits directly. It articulates what the quantitative approach can and cannot demonstrate, identifies the specific extensions that would strengthen each of the chapter's theoretical claims, and specifies a future research agenda---centered on direct qualitative inquiry into the drafting processes internal to legislative offices and advocacy organizations---that would test the mechanism-level propositions developed here. That agenda is not a concession that the current findings are incomplete. It is a recognition that different research instruments answer different questions, and that the textual and documentary evidence assembled in this thesis answers the questions it was designed to answer rigorously while bracketing the questions it was not designed to answer for subsequent investigation.
 
The chapter that follows develops these boundaries in detail.